\section*{Capítulo \ref{ch:g3:separating-mixtures} --- Separar misturas}

\begin{solution}{exo:g3:separating-mixtures:1}
\omterm{def:g3:separating-mixtures:sieve}{Peneiração}: os grãos de arroz são menores que as ervilhas e passam por
uma peneira com os furos certos. (Separar com a mão também funciona,
mas devagar.)
\end{solution}

\begin{solution}{exo:g3:separating-mixtures:2}
A terra desce devagar e forma uma camada no fundo; a água de cima fica
mais límpida. Isso é a \omterm{def:g3:separating-mixtures:decant}{sedimentação}.
\end{solution}

\begin{solution}{exo:g3:separating-mixtures:3}
O \omterm{def:g3:separating-mixtures:filter}{filtrado} é o líquido que passa pelo filtro. Quando se passa um café, o
\omterm{def:g3:separating-mixtures:filter}{filtrado} é o café na jarra; o pó fica no papel.
\end{solution}

\begin{solution}{exo:g3:separating-mixtures:4}
Sim. O sal está dissolvido, e um filtro não segura o que está
dissolvido: o sal passa junto com a água.
\end{solution}

\begin{solution}{exo:g3:separating-mixtures:5}
Fazendo a água evaporar: no sol, ou com um aquecimento suave, a água vai
para o ar e o açúcar fica no recipiente.
\end{solution}

\begin{solution}{exo:g3:separating-mixtures:6}
O cascalho fica na peneira. Os grãos de areia são menores que os furos,
então eles passam.
\end{solution}

\begin{solution}{exo:g3:separating-mixtures:7}
(a) \omterm{def:g3:separating-mixtures:sieve}{Peneiração}. (b) \omterm{def:g3:separating-mixtures:filter}{Filtração} (ou \omterm{def:g3:separating-mixtures:decant}{sedimentação} e \omterm{def:g3:separating-mixtures:decant}{decantação}). (c) \omterm{def:g3:separating-mixtures:evaporate}{Evaporação}.
\end{solution}

\begin{solution}{exo:g3:separating-mixtures:8}
Quatro etapas: grades, tanques de \omterm{def:g3:separating-mixtures:decant}{sedimentação}, filtros de areia,
eliminação dos micróbios. Os tanques de \omterm{def:g3:separating-mixtures:decant}{sedimentação} retiram a lama.
\end{solution}

\begin{solution}{exo:g3:separating-mixtures:9}
Filtrar; recolher a areia do papel; fazer o \omterm{def:g3:separating-mixtures:filter}{filtrado} evaporar para
obter o sal.
\end{solution}

\begin{solution}{exo:g3:separating-mixtures:10}
Não. O sal da água do mar está dissolvido, e um filtro não segura o que
está dissolvido: o \omterm{def:g3:separating-mixtures:filter}{filtrado} continua sendo água salgada.
\end{solution}

\begin{solution}{exo:g3:separating-mixtures:11}
Despejar o chá com cuidado depois que as folhas se depositaram
(\omterm{def:g3:separating-mixtures:decant}{decantação}), ou passá-lo por um coador ou por um filtro (\omterm{def:g3:separating-mixtures:sieve}{peneiração} ou
\omterm{def:g3:separating-mixtures:filter}{filtração}).
\end{solution}
